\section{The actual adjoint-defect source quotient and its continuous
dual}\label{the-actual-adjoint-defect-source-quotient-and-its-continuous-dual}

Independent review and receiving calculation. Proof locators
\textbf{AST0--AST9}.

\subsection{AST0. Construction stage, inputs and review
coverage}\label{ast0.-construction-stage-inputs-and-review-coverage}

The supporting datum remains
\(\tau\langle Z_1;\mathrm{no}\ Z_2\rangle\). All vector spaces,
arithmetic parameters, analytic coordinates and quotients below belong
to the already reconstructed coefficient system. The two branches retain
their separate recovered counters. No arithmetic, metric, midpoint or
vector operation is assigned to the support.

The connected reconstruction, its correction-precedence table, the
current construction guide, and verbatim WU062 and WU064--WU065 were
reread before this calculation. The new operation is a quotient of the
actual ACD specialization source, followed by continuous duality; its
prerequisites are proved below. The user's desired weight conclusion is
not added as a hypothesis.

The complete written proof
\href{https://github.com/KokunoYumeto/zeta-function-research-reader/blob/main/workbenches/splitzero-tandem/continuations/20260925-source-endpoint-and-section/GYSIN_DEFECT_CLOSURE_AND_POSITIVE_RECEIVER.md}{GDC0--GDC14}, final
SHA256
\nolinkurl{5d0f471b3250e4ee894caca68232c556d64a137a49a9fe25a641db61b94f79ab},
was independently read and checked in full. Its earlier version
\nolinkurl{f9e15dff7755b0d9c79c9dd896e55966b87302051fcf8086fad90e56da126589}
was also fully read; the final typography repair preserves the proof.
The review checked its reflection sign, exact annihilator, strong
closure, bounded parameter changes, all restriction and quotient
topologies, receiver triangle, positive-residue restriction, finite
multiplicities and source scope. A separate bounded independent check of
GDC7 and the invoked SDT4 compact-lift construction also found no
defect. This is a completed independent audit of the stated final
version, in addition to its earlier root verification.

The receiving proof uses the actual maps of
\href{https://github.com/KokunoYumeto/zeta-function-research-reader/blob/main/workbenches/splitzero-tandem/continuations/20260925-source-endpoint-and-section/CC_ADJOINT_DEFECT_CONTINUOUS_DUAL.md}{ACD0--ACD8}, SHA256
\nolinkurl{1373884c33e74aea9c0d48404152b16812b434b6cc4760227c8596666e919fda},
previously read and independently checked in full, and the explicit
seminorm, finite-net and compact-lift arguments of
\href{https://github.com/KokunoYumeto/zeta-function-research-reader/blob/96b02efaf4511d07f33de5ea4cfb10058e905d7e/workbenches/splitzero-tandem/continuations/20260924-cohomology-weight-and-character-lifts/gct-weight-control/proofs/CC_STRONG_DUAL_TOPOLOGY_AND_JETS.md}{SDT1--SDT4}, reread for this
extension. These are programme derivations. The human coefficient source
remains Connes--Consani, \emph{Schemes over} \(\mathbb F_1\) \emph{and
zeta functions},
\href{https://arxiv.org/abs/0903.2024v3}{arXiv:0903.2024v3, §5}; the
retained author TeX at
\nolinkurl{sources/CC_0903_2024_v3/author_source/announc3.tex},
lines1376--1667, was read in the preceding source reconstruction. No new
whole-paper reading or identification with Deligne's étale
invariant-cycle cross is claimed.

\subsection{AST1. The actual coefficient, the full original divisor and
the
defect}\label{ast1.-the-actual-coefficient-the-full-original-divisor-and-the-defect}

Retain the original spaces and maps \[
S=\{f\in\mathcal S(\mathbb R;\mathbb C):f(-v)=f(v),\ f(0)=0,
\ \int_{\mathbb R}f(v)\,dv=0\},
\] \[
A=\{a\in C^\infty(\mathbb R_{>0}):
p_{N,j}(a)=\sup_{u>0}(u^N+u^{-N})|(u\partial_u)^ja(u)|<\infty
\text{ for every }N,j\ge0\},
\] \[
\Sigma f(u)=2\sum_{k\ge1}f(ku),\quad J=\Sigma S,
\quad q:A\to Q=A/J,\quad B=A'_\beta.
\tag{AST1.1}
\] The accepted exact-image theorem makes \(J\) closed. The Mellin
comparison, with its original factor, is \[
\Theta a(s)=\frac12\int_0^\infty a(u)u^s\frac{du}{u},\qquad
\Theta\Sigma f(s)=\zeta(s)\int_0^\infty f(v)v^s\frac{dv}{v}
\quad(\Re s>1).
\tag{AST1.2}
\] Its entire rapid-strip quotient is \(Q\simeq\mathcal B/\mathcal I\),
where \(\mathcal I\) imposes all multiplicities of the actual nontrivial
zeros. The original source vector \[
f_0(v)=\frac\pi2v^2(2\pi v^2-3)e^{-\pi v^2}
\] has exactly \[
F_0(s)=\Theta\Sigma f_0(s)
=\frac{s(s-1)}8\pi^{-s/2}\Gamma(s/2)\zeta(s),
\] \[
F_0(0)=F_0(1)=\frac18,\quad F_0(-1)=F_0(2)=\frac\pi{24},\qquad
F_0(-2k)=\frac{k(2k+1)(-1)^k\pi^k}{2\,k!}\zeta'(-2k)
\quad(k\ge1).
\tag{AST1.3}
\] These values include the cancellations of the original pole and
trivial zeros; none creates an extra quotient jet. At an actual
nontrivial zero \(\rho\), writing \(z=s-\rho\), the retained local
expression is \[
F_0(\rho+z)=z^{m_\rho}
\frac{(\rho+z)(\rho+z-1)}8\pi^{-(\rho+z)/2}
\Gamma((\rho+z)/2)\frac{\zeta(\rho+z)}{z^{m_\rho}}.
\tag{AST1.4}
\] Every derivative in this product uses the full Leibniz sum.
Explicitly, away from its individually singular factors and then by
holomorphic continuation where the whole expression is regular, \[
F_0^{(j)}(s)=\sum_{k_0+k_1+k_2+k_3=j}
\frac{j!}{k_0!k_1!k_2!k_3!}
\left(\frac{s(s-1)}8\right)^{(k_0)}
\left(-\frac{\log\pi}{2}\right)^{k_1}\pi^{-s/2}
2^{-k_2}\Gamma^{(k_2)}(s/2)\zeta^{(k_3)}(s).
\tag{AST1.5}
\] This entire comparison remains a comparison with original \(\zeta\),
not its replacement. In its original half-plane the unit term, Euler
product and logarithmic derivative remain \[
\zeta(s)=1+\sum_{k\ge2}k^{-s}=\prod_p(1-p^{-s})^{-1},\qquad
-\frac{\zeta'(s)}{\zeta(s)}=\sum_p\sum_{k\ge1}(\log p)p^{-ks}.
\tag{AST1.6}
\]

Let \(\mathscr Z\) be the actual distinct nontrivial zeros, with
multiplicities \(m_\rho\), and set \(\rho^\#=1-\overline\rho\). Use \[
H=\ell^2(\mathscr Z,m),\quad
\langle x,y\rangle_+=\sum_\rho m_\rho x_\rho\overline{y_\rho},
\quad (\mathsf Jy)_\rho=y_{\rho^\#},\quad E:Q\to H.
\tag{AST1.7}
\] Here \((EF)_\rho=(\Theta_QF)(\rho)\), and the inner product is linear
in the first variable. The proved full-jet isolators realize every
finite value vector, so \(E\) is continuous with dense image. Put
\(H_L=\ell^2(\mathscr Z_L,m)\), \(H_O=\ell^2(\mathscr Z_O,m)\), where
the subsets denote respectively \(\Re\rho=1/2\) and \(\Re\rho\ne1/2\),
and write \(P_L,P_O\) for their orthogonal projections. These
definitions do not assert that either subset is empty or nonempty.

For a coefficient parameter \(t>0\), \((T_ty)_\rho=t^\rho y_\rho\) on
\(H\); on \(A\) the original action is \((T_ta)(u)=a(u/t)\), descending
to \(Q\) and giving the stated value action under \(E\). For the fixed
coefficient parameter \(r>1\), ACD uses \[
D_r=T_r^*-rT_{1/r},\quad
d_r(\rho)=e^{-i\Im\rho\log r}
(r^{\Re\rho}-r^{1-\Re\rho}),\quad
\beta_r=D_r^*\mathsf JE:Q\to H,\quad b_r=\beta_rq.
\tag{AST1.8}
\] Thus, with no multiplicity factor absorbed into the vector
coordinate, \[
(\beta_rF)_\rho=\overline{d_r(\rho)}F(\rho^\#).
\tag{AST1.9}
\] For \(C=\overline H\), the Riesz map is \[
R:C\to H',\quad R(\overline y)(h)=\langle h,y\rangle_+.
\] The exact continuous transposes are \[
\sigma_r=\beta_r'R:C\to Q'_\beta,\quad
a_r=b_r'R=q'\sigma_r:C\to B,
\] \[
\sigma_r(\overline y)(F)=\langle\beta_rF,y\rangle_+
=\langle EF,\mathsf JD_ry\rangle_+.
\tag{AST1.10}
\] Continuity into strong duals follows from the estimate on every
bounded test set \(K\), \(\sup_{F\in K}|\langle\beta_rF,y\rangle_+|
\le\sup_{F\in K}\|\beta_rF\|_+\|y\|_+\).

\subsection{AST2. The specialization kernel is exactly the GDC
source}\label{ast2.-the-specialization-kernel-is-exactly-the-gdc-source}

Because \(r>1\), the scalar \(d_r(\rho)\) vanishes exactly on
\(\mathscr Z_L\). Reflection preserves both subsets and their
multiplicities. Formula (AST1.9) therefore proves \[
\boxed{Z_r=\ker\beta_r=N_O:=\ker(P_OE),\qquad
\ker b_r=A_O:=q^{-1}N_O.}
\tag{AST2.1}
\] Both kernels are closed in their specified Fréchet spaces, and they
are independent of \(r>1\). The actual common specialization source is
consequently \[
\mathcal R:=Q/N_O\xleftarrow{\ \sim\ }A/A_O,
\qquad [a]\longmapsto[q(a)].
\tag{AST2.2}
\] This is a topological quotient isomorphism. Indeed \(q\) and the
quotient \(Q\to Q/N_O\) are open; their composite has kernel \(A_O\). An
open set in \(A/A_O\) is the image of an open set in \(A\), whose
composite image is open in \(Q/N_O\). The inverse is therefore
continuous, without selecting any section.

In the genuine ACD restriction fibre the quotient coordinate is
\((h,a)\mapsto h-b_ra\). Write \(\vartheta_p=d\arg z_p/(2\pi)\) for the
positive angular form in the oriented pole coordinate. Its boundaries
and their source comparison are exactly \[
\partial_{-1}=-b_r:A\to H,\qquad
\partial_0=+1:A[\vartheta_p]\to A,
\] \[
\operatorname{coker}\operatorname{sp}_{-1}
=A/A_O\simeq\mathcal R
\xrightarrow{\ -\overline\beta_r\ }b_rA,\qquad
\overline\beta_r([F])=\beta_rF.
\tag{AST2.3}
\] The last arrow is a continuous bijection onto the image with its
Hilbert subspace topology. It is a topological isomorphism if that image
is instead given its transported quotient topology. No inverse
continuity into the Hilbert subspace topology is asserted. The square
with the original boundary \(-q:A\to Q\) commutes because
\(b_r=\beta_rq\). Thus (AST2.2) is a refinement of the actual
specialization calculation, not a quotient selected independently of
that boundary.

At each actual multiplicity block \(\mathbb C[z]/z^{m_\rho}\), the
kernel \(N_O\) retains the whole block on \(\mathscr Z_L\), and retains
\(z\mathbb C[z]/z^{m_\rho}\) on \(\mathscr Z_O\). Accordingly
\(\mathcal R\) observes one off-line value coordinate and no line
coordinate on that block. All higher jets remain in the explicitly
retained kernel. This local description neither replaces the full global
quotient by a product nor proves unrestricted interpolation.

\subsection{AST3. The quotient topology and its full strong
dual}\label{ast3.-the-quotient-topology-and-its-full-strong-dual}

We prove the new quotient's bounded-compact property rather than
assuming a quotient theorem for general Montel spaces. Put \[
P_n(a)=\max_{0\le j\le n}\sup_{x\in\mathbb R}
(e^{nx}+e^{-nx})\left|\frac{d^j}{dx^j}a(e^x)\right|,
\quad
P_n^O([a])=\inf_{v\in A_O}P_n(a+v).
\tag{AST3.1}
\] The quotient topology on \(\mathcal R=A/A_O\) is precisely the
topology of these increasing seminorms. Closedness of \(A_O\) proves
Hausdorffness. Completeness follows by the SDT1 successive-lift argument
with \(A_O\) in place of \(J\): for a Cauchy subsequence choose lifts of
successive differences with \(P_k<2^{-k}\); their series converges in
\(A\) and maps to the required quotient limit.

Let \(K\subset\mathcal R\) be bounded. For each fixed \(n\), choose
representatives of its elements with \[
P_{n+1}(a_y)<1+\sup_{z\in K}P_{n+1}^O(z).
\tag{AST3.2}
\] SDT2 proves that the set with this single \(P_{n+1}\) bound is
totally bounded for \(P_n\), using both original exponential weights and
one additional derivative. Its finite net maps to a finite \(P_n^O\)-net
for \(K\). This holds for every \(n\), so \(K\) is totally bounded in
the complete quotient metric. Its closure is compact. Thus
\(\mathcal R\) is a Montel Fréchet space by this direct calculation.

The strict quotient \(\pi:Q\to\mathcal R\) therefore lifts every bounded
set to a bounded set in \(Q\). Indeed take the compact closure just
proved, apply SDT4's finite-choice compact-lift construction to this
strict Fréchet quotient, and intersect the resulting compact lift with
the preimage of the original bounded set. The image of that bounded lift
is exactly the original set. The same construction applies to
\(\widetilde\pi=\pi q:A\to\mathcal R\).

Continuous descent of functionals gives algebraic bijections onto the
indicated annihilators. Their strong topologies agree because for any
bounded \(K\subset\mathcal R\) and an exact bounded lift
\(\widetilde K\subset Q\), \[
\sup_{x\in K}|\lambda(x)|
=\sup_{F\in\widetilde K}|(\pi'\lambda)(F)|.
\tag{AST3.3}
\] The forward seminorm estimate follows by mapping any bounded set of
\(Q\) to a bounded set of \(\mathcal R\). Thus \[
\boxed{\pi':\mathcal R'_\beta\xrightarrow{\sim}S_O=N_O^\perp
\subset Q'_\beta,\qquad
\widetilde\pi':\mathcal R'_\beta\xrightarrow{\sim}A_O^\perp
\subset B.}
\tag{AST3.4}
\] These are closed strong topological embeddings, and
\(\widetilde\pi'=q'\pi'\). This proves the exact dual of the actual ACD
source quotient.

GDC7's independently checked restriction theorem now gives the strict
exact rows \[
0\to\mathcal R'_\beta\xrightarrow{\pi'}Q'_\beta
\xrightarrow{\mathrm{res}}(N_O)'_\beta\to0,
\] \[
0\to\mathcal R'_\beta\xrightarrow{\widetilde\pi'}B
\xrightarrow{\mathrm{res}}(A_O)'_\beta\to0.
\tag{AST3.5}
\] The restrictions are strongly open, not merely onto as vector maps.
Their bounded dual families also have bounded lifts by the common
Hahn--Banach seminorm estimate proved in GDC7. The embeddings, openness
and bounded lifts used here have separate proofs; no continuous linear
section is inferred.

\subsection{AST4. The actual pairing is strongly dense in this
dual}\label{ast4.-the-actual-pairing-is-strongly-dense-in-this-dual}

The map \(\overline\beta_r:\mathcal R\to H_O\) is continuous and
injective by (AST2.1). Its image is dense in \(H_O\): for a finite
vector \(h\) on off-line zeros, prescribe the finite values
\(F(\rho^\#)=h_\rho/\overline{d_r(\rho)}\), with zero values elsewhere,
using the already proved full-jet isolators. Formula (AST1.9) then gives
exactly \(\beta_rF=h\). This proves Hilbert density, without claiming
density of finite blocks in \(Q\).

Define the continuous linear transpose receiver \[
\gamma_r:\overline{H_O}\to\mathcal R'_\beta,
\qquad
\gamma_r(\overline y)([F])=\langle\beta_rF,y\rangle_+.
\tag{AST4.1}
\] It is injective by that density and has the exact comparison \[
\pi'\gamma_r=\sigma_r|_{\overline{H_O}},\qquad
\widetilde\pi'\gamma_r=a_r|_{\overline{H_O}}.
\tag{AST4.2}
\] The strong topological isomorphism in (AST3.4), together with GDC2's
proved strong closure, gives \[
\overline{\gamma_r(\overline{H_O})}^{\,\beta}=\mathcal R'_\beta.
\tag{AST4.3}
\] Thus the nondegenerate ACD pairing on
\(\mathcal R\times\overline{H_O}\) now has its exact dual receiver. It
need not exhaust that dual before closure.

Its Hilbert norm on the source is exactly \[
\|[F]\|_r^2:=\|\overline\beta_r[F]\|_+^2
=\sum_{\rho\in\mathscr Z_O}m_\rho|d_r(\rho)|^2|F(\rho^\#)|^2
=\sum_{\rho\in\mathscr Z_O}m_\rho|d_r(\rho)|^2|F(\rho)|^2.
\tag{AST4.4}
\] The second equality reindexes by \(\#\), using equal multiplicities
and \(d_r(\rho^\#)=-d_r(\rho)\). Completion for this displayed norm is
\(H_O\), via \(\overline\beta_r\). This norm topology is not equated
with the original Fréchet quotient topology.

\subsection{AST5. The split dual receiver is the transpose of an actual
primal
subcomplex}\label{ast5.-the-split-dual-receiver-is-the-transpose-of-an-actual-primal-subcomplex}

At a pole, ACD's actual primal object is \[
\mathscr P_r=[\underline A^{-1}\xrightarrow{+b_r}i_*H^0].
\] Its source kernel constructs the literal subcomplex \[
\mathscr P_O=[\underline{A_O}^{-1}\xrightarrow{0}i_*H^0]
=\underline{A_O}[1]\oplus i_*H.
\tag{AST5.1}
\] The map \(\mathscr P_O\to\mathscr P_r\) is inclusion on \(A_O\) and
identity on \(H\). It is a chain map because \(b_r|_{A_O}=0\). Its
literal quotient is the constant sheaf \(\underline{\mathcal R}[1]\), so
there is a short exact sequence of the actual complexes \[
0\to\mathscr P_O\longrightarrow\mathscr P_r
\xrightarrow{(\widetilde\pi,0)}\underline{\mathcal R}[1]\to0.
\tag{AST5.2}
\] Each coefficient inclusion/quotient has its proved original topology.
The quotient complex itself has zero differential. First define the
ordinary short-exact-sequence boundary by differentiating a lifted
cocycle with a positive sign. A stalk class
\([a]\in H^{-1}(\underline{\mathcal R}[1])=\mathcal R\) lifts to
\(a\in A=\mathscr P_r^{-1}\); its differential is
\(+b_ra\in H=\mathscr P_O^0\). Changing the lift by \(v\in A_O\) changes
this differential by \(b_rv=0\). Hence \[
\delta_{\rm SES}([a])=+b_ra=+\overline\beta_r([a]).
\tag{AST5.2a}
\] There is a separate sign in the retained mapping-cone convention. For
the inclusion \(i:\mathscr P_O\to\mathscr P_r\), take
\(\operatorname{Cone}(i)^k=\mathscr P_r^k\oplus\mathscr P_O^{k+1}\),
\(d(v,u)=(dv+i u,-du)\), and last cone arrow given by positive
projection to \(\mathscr P_O[1]\). The quasi-isomorphism from this cone
to \(\underline{\mathcal R}[1]\) is the positive quotient on its
\(\mathscr P_r\) coordinate. At the pole the cone has \(A\oplus H\) in
degree \(-1\), with differential \((a,h)\mapsto b_ra+h\). The cocycle
\((a,-b_ra)\) represents \([a]\), and its positive last projection is
\(-b_ra\). Consequently the distinguished triangle with the literal
positive quotient is \[
\mathscr P_O\xrightarrow{i}\mathscr P_r
\xrightarrow{+\widetilde\pi}\underline{\mathcal R}[1]
\xrightarrow{\kappa_r}\mathscr P_O[1],\qquad
H^{-1}(i^*\kappa_r)=-\overline\beta_r=-\delta_{\rm SES}.
\tag{AST5.2b}
\] Changing the quotient coordinate by \(-1_{\mathcal R}\) instead gives
the isomorphic triangle with middle arrow \(-\widetilde\pi\), last arrow
\(-\kappa_r\), and positive stalk boundary \(+\overline\beta_r\). Thus
the positive differentiate-a-lift boundary and the stated cone
projection have their exact coordinate comparison; their signs have not
been identified silently.

Both signs of this map are continuous by the original quotient topology,
injective, and dense in \(H_O\) by AST4. A lift of an individual class
is all the calculation requires; it supplies no continuous section of
\(A\to\mathcal R\). The actual localization boundary in (AST2.3) is
\(-b_r\) because its restriction-fibre quotient coordinate is
\(h-b_ra\). It agrees with \(H^{-1}(i^*\kappa_r)\widetilde\pi\) as a
coefficient map, while the positive short-exact-sequence convention
gives the negative of that map. These are distinct constructions with
the stated exact relation.

Use ACD's fine resolution with the minus de Rham differential. The
inclusion of smooth \(A_O\)-valued forms into smooth \(A\)-valued forms
and the identity on the point coefficient give a continuous
fine-resolution map. Its literal compact-test transpose, with ACD's
degree signs \((-1,+1,-1)\) and Riesz identification, is \[
\mathscr E_r\longrightarrow
\mathscr T_{A_O}[-1]\oplus i_*\overline H,
\qquad (u,x)\longmapsto(\iota_O^\vee u,x).
\tag{AST5.3}
\] Here \(\mathscr E_r=\operatorname{Cone}(-\delta_pa_r)[-1]\). The
target attaching map is zero because the coefficient restriction
\(B\to(A_O)'\) kills \(a_r\). Explicitly
\(\iota_O^\vee\delta_pa_r=\delta_p(a_r|_{A_O})=0\); hence the current
differential commutes. This is exactly GDC10.3, now identified as the
continuous transpose of (AST5.2)'s first map.

The compact-test contraction applies to the complete Fréchet coefficient
\(A_O\), giving \(\mathscr T_{A_O}[-1]\simeq\underline{(A_O)'}[1]\). On
underlying vector-space sheaves, (AST3.5) and cancellation of the common
supported term give the exact receiver triangle \[
\boxed{\underline{\mathcal R'}[1]\longrightarrow\mathscr E_r
\longrightarrow\underline{(A_O)'}[1]\oplus i_*\overline H
\longrightarrow\underline{\mathcal R'}[2].}
\tag{AST5.4}
\] Under \(\widetilde\pi'\), its first coefficient is exactly
\(A_O^\perp\), and the triangle is GDC10.4. For a direct verification of
its fibre, form the two current Gysin triangles with common supported
coordinate. Their comparison cancels that coordinate; compact-test
contractions reduce the remaining comparison to the actual coefficient
restriction \(B\to(A_O)'\), whose kernel is \(\mathcal R'\) by (AST3.5).
The fibre is thus \(\underline{\mathcal R'}[1]\). No general exactness
claim for strong duals of arbitrary complexes is used.

The positive costalk map is consequently \[
[\overline H\xrightarrow{+a_r}B]
\longrightarrow[\overline H\xrightarrow{0}(A_O)'],
\qquad(x,b)\longmapsto(x,b|_{A_O}).
\tag{AST5.5}
\] Its degree-one map is the separated quotient of the actual
unseparated cohomology: \[
B/a_rC\longrightarrow B/A_O^\perp\xrightarrow{\sim}(A_O)'_\beta,
\quad
\ker= A_O^\perp/a_rC.
\tag{AST5.6}
\] That kernel is precisely the closure of zero in the original costalk
quotient. The retained source \(\mathcal R\) and its dual kernel in
(AST5.4) identify the exact cost of this split receiver.

\subsection{AST6. Parameter isomorphisms are actual transpose
comparisons}\label{ast6.-parameter-isomorphisms-are-actual-transpose-comparisons}

For \(r,t>1\), retain GDC3's bounded invertible diagonal map \[
v_{r,t}(\rho)=
\begin{cases}
e^{-i\Im\rho\log(r/t)}
\dfrac{r^{\Re\rho}-r^{1-\Re\rho}}
{t^{\Re\rho}-t^{1-\Re\rho}},&\rho\in\mathscr Z_O,\\
1,&\rho\in\mathscr Z_L,
\end{cases}
\qquad V_{r,t}y=(v_{r,t}(\rho)y_\rho)_\rho.
\tag{AST6.1}
\] Its explicit off-line modulus bounds are \[
\frac{2\sqrt r\log r}{(t+1)\log t}
\le |v_{r,t}(\rho)|
\le\frac{(r+1)\log r}{2\sqrt t\log t};
\tag{AST6.2}
\] on line coordinates its value is one. The nonzero constants prove
bounded invertibility, even for zeros approaching the line. Put
\(M_{r,t}=V_{r,t}^*\). Since \(D_r=D_tV_{r,t}\), taking Hilbert adjoints
gives the exact identities \[
\beta_r=M_{r,t}\beta_t,\qquad
b_r=M_{r,t}b_t,
\qquad \sigma_t\overline V_{r,t}=\sigma_r,
\qquad a_t\overline V_{r,t}=a_r.
\tag{AST6.3}
\] In particular \[
\mathscr P_t\xrightarrow{\ (1_A,M_{r,t})\ }\mathscr P_r,
\qquad
\mathscr E_r\xrightarrow{\ (1_B,\overline V_{r,t})\ }\mathscr E_t
\tag{AST6.4}
\] are continuous cochain isomorphisms with inverses obtained by
exchanging \(r,t\). The first has the same sign on its positive primal
differential; the second has the same sign on its positive Gysin
attaching term. They are literal continuous transposes under ACD's
pairing, because \[
M_{r,t}'R(\overline y)(h)
=\langle M_{r,t}h,y\rangle_+
=\langle h,V_{r,t}y\rangle_+
=R(\overline V_{r,t}\overline y)(h).
\tag{AST6.5}
\] The diagonal operators commute and satisfy the GDC cocycle identity,
so these maps satisfy the corresponding compositional identities in
their indicated directions. They fix \(A,A_O,Q,N_O,\mathcal R\), the
entire nearby dual coefficient \(B\), and the line-supported coefficient
\(H_L\) or \(\overline{H_L}\). They do not identify the numerical defect
sizes for different parameters.

The localization boundary squares are now explicit: \[
-b_r=M_{r,t}(-b_t),\qquad
1_A=1_A\,1_A,\qquad
a_r=a_t\overline V_{r,t}.
\tag{AST6.6}
\] The induced map on the quotient
\(\operatorname{coker}\operatorname{sp}_{-1}=\mathcal R\) is identity;
its map between the displayed Hilbert images is \(M_{r,t}\). Its map on
the dual source \(\mathcal R'\) in (AST5.4) is identity, and
\(\gamma_t\overline V_{r,t}=\gamma_r\) on \(\overline{H_O}\).

These parameter maps preserve the actual cover structures. For a
recovered integer degree \(n\ge1\), \(V_{r,t}\) and \(M_{r,t}\) commute
with every diagonal \(T_n,U_n\) and their Hilbert adjoints. Thus
(AST6.4) commutes with the primal unit and weighted trace and the dual
DCA maps, including their finite-sheet corrections. In DCA's correction
vector every occurrence of \(R_-a_r x\) equals
\(R_-a_t\overline V_{r,t}x\), so all sheet coordinates commute, not
merely their sums. The integer cover and its angular degree factors are
retained; no geometric cover of noninteger degree is asserted.

\subsection{AST7. Both poles, the global specialization map and
endpoints}\label{ast7.-both-poles-the-global-specialization-map-and-endpoints}

At both poles use the same existing coefficient trivialization. The
subcomplex construction becomes \[
\mathscr P_O^Y=\underline{A_O}[1]\oplus i_+H\oplus i_-H
\longrightarrow
\mathscr P_r^Y=[\underline A^{-1}\xrightarrow{(b_r,b_r)}i_+H\oplus i_-H],
\tag{AST7.1}
\] with quotient \(\underline{\mathcal R}_Y[1]\). The source and dual
maps are identity on each original endpoint summand \(E_p[1]\) and
\(E_p'[-1]\). Both distinct value-at-zero and integral labels remain;
none enters the kernel \(N_O\) by a relabelling of a zeta endpoint.

In ACD6's proved global finite model the parameter map is \[
[A^{-1}\xrightarrow{(b_t,b_t)}H^0\oplus H^0\xrightarrow0A^1]
\longrightarrow
[A^{-1}\xrightarrow{(b_r,b_r)}H^0\oplus H^0\xrightarrow0A^1],
\] \[
(1_A,M_{r,t}\oplus M_{r,t},1_A).
\tag{AST7.2}
\] It induces identity on \(A_O\oplus E_+\oplus E_-\) in degree \(-1\),
identity on \(A\) in degree one, and \[
H\oplus(H/b_tA)\longrightarrow H\oplus(H/b_rA),
\quad(h,[k])\longmapsto(M_{r,t}h,[M_{r,t}k])
\tag{AST7.3}
\] in degree zero. This is well defined with a continuous inverse
because \(M_{r,t}b_tA=b_rA\), and quotient topology gives continuity
even if these quotients are unseparated. The original global source map
remains \[
Q\longrightarrow H\oplus(H/b_rA),\qquad
F\longmapsto(\beta_rF,0),\qquad \ker=N_O.
\tag{AST7.4}
\] Equations (AST6.3) and (AST7.3) prove its parameter compatibility on
all original classes.

The dual finite global model has \(B^{-1}\), \(C^0\oplus C^0\), \(B^1\)
and differential \((x_+,x_-)\mapsto a_rx_++a_rx_-\). Its parameter map
is identity on both \(B\) terms and \(\overline V_{r,t}\) on each point
term. The two positive Dirac contributions remain a sum. Since GDC3
proves equality of the actual images \(a_rC=a_tC\), the degree-one
quotient maps are literally identity on \(B/a_rC\), and their separated
maps are identity on \((A_O)'\). Endpoints extend these statements by
their identity maps. This proves preservation of the global receiving
maps without declaring an anti-diagonal cancellation to be vanishing of
the local attaching map.

\subsection{AST8. The restricted original extension and the surviving
positive
residue}\label{ast8.-the-restricted-original-extension-and-the-surviving-positive-residue}

The actual kernel source belongs to the strict original pullback row \[
0\to J\to A_O\xrightarrow{q_O}N_O\to0.
\tag{AST8.1}
\] Its continuous strong dual is, by GDC7, \[
0\to(N_O)'_\beta\xrightarrow{q_O'}(A_O)'_\beta
\to J'_\beta\to0.
\tag{AST8.2}
\] This row retains all of \(J\), in contrast to replacing the
coefficient by a zeta value space. It also proves that the split
receiver's nearby coefficient still contains the dual of the source
carrying all line values and all higher jets.

For completeness, its retained original residue is the proved exact
restriction \[
A_H(y)|_{N_O}(F)
=\sum_{\rho\in\mathscr Z_L}m_\rho F(\rho)\overline{y_\rho}
=\langle EF,P_Ly\rangle_+\qquad(F\in N_O).
\tag{AST8.3}
\] It injects \(\overline{H_L}\) into \((N_O)'\), by testing each actual
finite line-value isolator, and then into \((A_O)'\) through \(q_O'\).
Precisely, \(N_O/N_0\), where \(N_0=\ker E\), with norm
\(\|[F]\|=\|EF\|_+\), completes to \(H_L\) by GDC8.3. The injected
\(H_L\) already is complete in its own Hilbert norm. Neither that norm
nor that Hilbert image is identified with the full strong dual
\((N_O)'\); the original higher-jet functionals remain present. The full
source remains the exact extension \(0\to N_O\to Q\to\mathcal R\to0\);
no source projector splitting this row has been constructed or assumed.

\subsection{AST9. Exact result and limits of the
comparison}\label{ast9.-exact-result-and-limits-of-the-comparison}

The actual ACD specialization source has now been identified as
\(\mathcal R=Q/N_O\), with its original quotient topology. Its full
strong dual is the closed GDC defect receiver \(S_O\), and its image in
the original dual coefficient is \(A_O^\perp\). The original
specialization map, its sign, its two-pole map and every parameter
comparison commute with these identifications. The GDC split current
receiver is the literal continuous transpose of the actual subcomplex
(AST5.1); the discarded coefficient is retained as
\(\underline{\mathcal R'}[1]\) in its exact triangle.

No existence or nonexistence of an off-line zero has entered this proof.
If an actual off-line zero is present, its full-jet isolator makes its
value class nonzero in \(\mathcal R\); if none is present then \(N_O=Q\)
and \(\mathcal R=0\). These are exact statements about the proved
receiver. They do not prove which case holds, impose a common weight, or
turn the coefficient quotient into Deligne's geometric inertia cross.
What advances the comparison is the exhibited primal subcomplex, its
quotient topology, its complete dual and its commuting specialization
maps, all on the original arithmetic coefficients.
